\subsection{Problem Motivation}

Fine-grained classification means telling apart visually similar categories: bird species, dog breeds, flower types. The model has to pick up on subtle differences in shape and colour while shrugging off variations in lighting and pose.

Oxford Flowers 102 \citep{nilsback2008} makes this harder by giving us very little to work with. The standard split has only 10 training images per class across 102 classes, which is barely enough to fit a softmax head, let alone an 11M-parameter network. That makes it a useful benchmark for few-shot methods.

CNNs and Vision Transformers (ViTs) do well on large datasets, but it is less obvious how they behave when training data is scarce. We wanted to find out which adaptation strategy holds up best at $k \in \{1, 2, 5, 10\}$ examples per class.

\subsection{Approach Overview}

We compare three families of adaptation under a single shared training recipe (AdamW, cosine annealing, 50 epochs, batch size 32) so that any difference in performance is attributable to the method itself:

\begin{itemize}
  \item \textbf{Five ResNet18 variants}: a baseline plus four advanced convolutional modifications (dilated, deformable, depthwise-separable, and our novel hybrid combining dilated and deformable layers).
  \item \textbf{Four regularization strategies} on top of the ResNet18 baseline: freezing early layers, label smoothing, MixUp, and stronger image augmentation.
  \item \textbf{A larger CNN backbone} (ResNet50) to test whether more capacity helps.
  \item \textbf{Visual Prompt Tuning (VPT)} on a ViT-B/16 backbone, in both Shallow and Deep configurations \citep{jia2022vpt}.
\end{itemize}

We measure accuracy, mean per-class accuracy, and top-5 accuracy on the test set. We also rerun every method at $k = 1, 2, 5$ images per class to see how it degrades. Two questions drive the work: does architectural complexity actually help in the few-shot regime, and can parameter-efficient tuning match or beat full fine-tuning on a fine-grained task?

\subsection{Training and Evaluation Setup}

We use the canonical \texttt{torchvision.datasets.Flowers102} split: 1020 training images, 1020 validation images, and 6149 test images. Images are resized to $224 \times 224$ and normalised with ImageNet statistics. The training transform applies a random horizontal flip and rotation up to 15 degrees.

We report three test-set metrics:
\begin{itemize}
  \item \textbf{Accuracy}: overall correct predictions over all samples.
  \item \textbf{Mean per-class accuracy (MPC)}: average per-class recall, our primary metric following \citet{nilsback2008}. The test set is class-imbalanced (20 to 258 images per class), so this corrects for class frequency.
  \item \textbf{Top-5 accuracy}: fraction of samples where the true class appears in the top-5 predictions.
\end{itemize}

Unless stated otherwise, every model trains for 50 epochs with AdamW ($\text{lr} = 10^{-3}$), cosine annealing, batch size 32, automatic mixed precision, and seed 69. Experiments run on a single NVIDIA RTX 3070 Ti.
